\documentclass[11pt]{article}

\usepackage[utf8]{inputenc}
\usepackage[T1]{fontenc}
\usepackage{lmodern}
\usepackage{geometry}
\usepackage{amsmath,amssymb,amsfonts,amsthm,mathtools}
\usepackage{bm}
\usepackage{enumitem}
\usepackage{microtype}
\usepackage{hyperref}
\usepackage{titlesec}
\usepackage{booktabs}
\usepackage{array}
\usepackage{setspace}
\usepackage{mathrsfs}
\usepackage{physics}

\geometry{margin=1in}
\hypersetup{colorlinks=true, linkcolor=black, urlcolor=blue, citecolor=black}
\setlist[enumerate]{topsep=0.4em,itemsep=0.35em}
\setlist[itemize]{topsep=0.3em,itemsep=0.25em}
\titleformat{\section}{\large\bfseries}{\thesection.}{0.5em}{}
\titleformat{\subsection}{\normalsize\bfseries}{\thesubsection.}{0.5em}{}
\setstretch{1.06}

\newcommand{\Aether}{\AE{}ther}
\newcommand{\Aflow}{\AE{}ther-flow}
\newcommand{\TheoryName}{\Aflow{} Interpretation of Relativity}
\newcommand{\TheoryProgram}{\Aether{} / \Aflow{} framework}

\newcommand{\CompletedMetric}{g^{\sharp}}

\newcommand{\Car}{\mathrm{car}}
\newcommand{\Cur}{\mathrm{cur}}
\newcommand{\Hom}{\mathrm{hom}}

\newcommand{\ForSpace}[1]{\mathcal I_{#1}^{\mathrm{for}}}
\newcommand{\PrimShell}{\mathcal S_{\mathrm{prim}}}

\newcommand{\Reduction}[1]{\mathcal R_{#1}}
\newcommand{\CurrentCarrier}[1]{\mathfrak C_{#1}^{\mathrm{cur}}}
\newcommand{\EqCarrier}[1]{\mathfrak C_{#1}^{\mathrm{eq}}}
\newcommand{\CurrentExtract}[1]{\mathscr E_{#1}^{\mathrm{cur}}}
\newcommand{\EqExtract}[1]{\mathscr E_{#1}^{\mathrm{eq}}}
\newcommand{\PrimCurrent}[1]{\mathcal J_{#1}^{\mathrm{prim}}}
\newcommand{\PrimTheta}[1]{\Theta_{#1}^{\mathrm{prim}}}

\newcommand{\MetricOut}{\mathscr G_{\mathrm{out}}}
\newcommand{\MatterOut}{\mathscr T_{\mathrm{out}}}

\newcommand{\VisibleAffine}[1]{\widehat X_{#1}}
\newcommand{\DataSpace}[1]{\mathcal Y_{#1}^{\mathrm{obs}}}
\newcommand{\AmbientTarget}[1]{E_{#1}^{\mathrm{amb}}}

\newcommand{\EqnSpace}[1]{\mathcal U_{#1}^{\mathrm{eqn}}}
\newcommand{\EqnBody}[1]{\mathcal B_{#1}^{\mathrm{eqn}}}
\newcommand{\EqnData}[1]{\Lambda_{#1}^{\mathrm{eqn,dat}}}
\newcommand{\EqnShell}[1]{\Lambda_{#1}^{\mathrm{eqn,sh}}}
\newcommand{\EqnTarget}[1]{Q_{#1}^{\mathrm{eqn}}}
\newcommand{\EqnPair}[1]{\mathfrak b_{#1}^{\mathrm{eqn}}}
\newcommand{\EqnResidual}[1]{\mathcal E_{#1}^{\mathrm{eqn}}}
\newcommand{\EqnZero}[1]{V_{#1}^{\mathrm{eqn},0}}
\newcommand{\EqnHidden}[1]{H_{#1}^{\mathrm{eqn}}}
\newcommand{\EqnProj}[1]{\Pi_{#1}^{\mathrm{eqn,hid}}}
\newcommand{\EqnLin}[1]{\mathcal N_{#1}^{\mathrm{eqn}}}
\newcommand{\EqnEmbed}[1]{\zeta_{#1}^{\mathrm{eqn}}}

\newcommand{\ResSpace}[1]{\mathcal U_{#1}^{\mathrm{sup,res}}}
\newcommand{\ResBody}[1]{\mathcal B_{#1}^{\mathrm{sup,res}}}
\newcommand{\ResTarget}[1]{S_{#1}^{\mathrm{sup,res}}}
\newcommand{\ResPair}[1]{\mathfrak m_{#1}^{\mathrm{sup,res}}}
\newcommand{\ResSection}[1]{\Theta_{#1}^{\mathrm{sup,res}}}
\newcommand{\ResData}[1]{\Lambda_{#1}^{\mathrm{sup,dat}}}
\newcommand{\ResShellMap}[1]{\Lambda_{#1}^{\mathrm{sup,sh}}}
\newcommand{\ResShellSet}[1]{\mathcal Z_{#1}^{\mathrm{sup,res}}}
\newcommand{\ResZeroCore}[1]{\mathcal Z_{#1}^{0,\mathrm{sup}}}
\newcommand{\SeedHidden}[1]{\mathcal H_{#1}^{\mathrm{sup}}}
\newcommand{\SeedHiddenBody}[1]{D_{#1}^{\mathrm{sup}}}
\newcommand{\ShellChart}[1]{\Psi_{#1}^{\mathrm{sup}}}
\newcommand{\ZeroChart}[1]{\Psi_{#1}^{0}}
\newcommand{\HiddenChart}[1]{\Psi_{#1}^{\mathrm{hid}}}
\newcommand{\SplitMap}[1]{\Phi_{#1}^{\mathrm{sup}}}
\newcommand{\ZeroBody}[1]{\mathcal B_{#1}^{0,\mathrm{sup}}}
\newcommand{\ResResidual}[1]{\mathcal R_{#1}^{\mathrm{sup}}}
\newcommand{\HiddenCharge}[1]{\mathcal Q_{#1}^{\mathrm{sup}}}
\newcommand{\ResEmbed}[1]{\zeta_{#1}^{\mathrm{sup}}}
\newcommand{\Kernel}[1]{K_{#1}^{0}}
\newcommand{\StateComp}[1]{P_{#1}^{\mathrm{co}}}

\newcommand{\SubSpace}[1]{\mathcal M_{#1}^{\mathrm{sub}}}
\newcommand{\SubBody}[1]{\mathcal B_{#1}^{\mathrm{sub}}}
\newcommand{\SubTarget}[1]{E_{#1}^{\mathrm{sub}}}
\newcommand{\SubPair}[1]{\mathfrak s_{#1}^{\mathrm{sub}}}
\newcommand{\SubSection}[1]{\Sigma_{#1}^{\mathrm{sub}}}
\newcommand{\SubFunctional}[1]{\mathscr U_{#1}^{\mathrm{sub}}}
\newcommand{\SubShellSet}[1]{\mathcal Z_{#1}^{\mathrm{sub}}}
\newcommand{\SubData}[1]{\Lambda_{#1}^{\mathrm{sub,dat}}}
\newcommand{\SubShellMap}[1]{\Lambda_{#1}^{\mathrm{sub,sh}}}
\newcommand{\SubSym}[1]{\mathbb G_{#1}^{\mathrm{sub}}}
\newcommand{\SubTime}[1]{\vartheta_{#1}^{\mathrm{sub}}}
\newcommand{\SubChart}[1]{\Xi_{#1}^{\mathrm{sub}}}
\newcommand{\SubSupport}[1]{\Theta_{#1}^{\mathrm{sub,sup}}}
\newcommand{\SubSupportShell}[1]{\mathcal Z_{#1}^{\mathrm{sub,sup}}}
\newcommand{\SubZeroCore}[1]{\mathcal Z_{#1}^{0,\mathrm{sub}}}
\newcommand{\SubZeroChart}[1]{\Psi_{#1}^{0,\mathrm{sub}}}
\newcommand{\SubSplit}[1]{\Phi_{#1}^{\mathrm{sub}}}
\newcommand{\SubCharge}[1]{\mathcal Q_{#1}^{\mathrm{sub}}}
\newcommand{\SubResidual}[1]{\mathcal R_{#1}^{\mathrm{sub}}}
\newcommand{\SubEmbed}[1]{\zeta_{#1}^{\mathrm{sub}}}
\newcommand{\SubKernel}[1]{K_{#1}^{0,\mathrm{sub}}}
\newcommand{\SubKernelCh}[1]{\widetilde K_{#1}^{0,\mathrm{sub}}}
\newcommand{\SubStateComp}[1]{P_{#1}^{\mathrm{sub,co}}}

\newtheorem{definition}{Definition}
\newtheorem{lemma}{Lemma}
\newtheorem{proposition}{Proposition}
\newtheorem{remark}{Remark}
\newtheorem{corollary}{Corollary}
\newtheorem{theorem}{Theorem}

\title{\textbf{The \TheoryName{}}\\[0.4em]
\large Primitive-Reservoir Observer-Reduction\\
\large Hidden-Translation Support Reservoir Dynamics\\
\large from Primitive Substrate Shell Restriction}
\author{}
\date{}

\begin{document}

\maketitle

\begin{abstract}
The new hidden-translation support-reservoir theorem on the active primitive-reservoir observer-reduction line sharpened the project burden one step further. The designated zero-hidden equation-support package and hidden-fiber exactness are no longer the honest stopping point there; they are forced once one grants the hidden-translation support reservoir package. The remaining burden recorded by that theorem is deeper again: explain why the reservoir package itself exists, or remove the remaining reservoir-level freedom by a sharper inevitability result.

The present note takes the origin-side route. For each sector it introduces a more primitive substrate shell-restriction dynamics: a compact convex substrate body carrying an exact restriction section with normal shell stability, affine observer-data and shell-response readouts, symmetry and time reversal, an exact shell chart whose image is the reservoir body, a shell-internal support recorder whose chart transport is the reservoir support section, an exact support shell with zero-core / hidden-seed decomposition, a support-shell chart and hidden charge, an observer-compatible zero-response realization, fixed-data solvability on the support zero core, and a transported inert-kernel coercivity inequality. The recorded hidden-translation support reservoir package is then recovered unchanged as the shell-chart image of that deeper substrate package rather than inserted independently.

Because the same reservoir package is recovered, the same downstream consequences survive unchanged: the designated support-layer data and hidden-fiber exact residual are forced again, with canonical residual
\[
\EqnResidual{\sigma}(z+\eta)=\EqnLin{\sigma}(\eta),
\]
the one operative metric remains exactly \(\CompletedMetric\), and the ordinary matter route is unchanged. The scope stays disciplined. This is an origin-side gain in the active derivational program of \TheoryName{}, not a license for stronger first-principles promotion language. The next honest burden is therefore lower still: derive or justify the present substrate shell-restriction package from yet more primitive substrate dynamics, or else prove a sharper inevitability theorem at that level.
\end{abstract}

\section{Introduction}

The active derivational program of \TheoryName{} has now reached a cleaner origin-side bottleneck than another support-level or observer-level reformulation. The immediately preceding theorem on the primitive-reservoir observer-reduction line already proved that the designated zero-hidden equation-support package and hidden-fiber exactness are forced once one grants the hidden-translation support reservoir package. That theorem matters precisely because it removes the temptation to stop at the support layer. The next honest burden is not another stationary-reduction restatement, not another stationary-Hessian or spectral repackaging, not another selector-level relay, and not any change to the one operative metric or the ordinary matter law. The next honest burden is to explain why the hidden-translation support reservoir package itself should hold.

That is the role of the present note. It starts below the reservoir package rather than beside it. The positive claim is that there is a still deeper substrate shell-restriction dynamics whose exact shell chart produces the reservoir body, the exact support section and its normal stability, the zero-core / hidden-seed decomposition, the shell chart, the hidden charge, the compatibility realization, the solvability statement, the inert kernel/projector, and the coercivity inequality. In that sense the present note is stronger than another same-output restatement and narrower than any claim of completed final microphysical closure.

The point of the manuscript is therefore disciplined and specific. It does not promote the project more broadly. It does not alter the benchmark exact-relativistic package. It only pushes the remaining honest burden one layer lower on the origin side by showing that the recorded hidden-translation support reservoir package is itself a transported exact-shell consequence of a more primitive substrate dynamics.

\paragraph{Framework and claim boundary.}
\TheoryProgram{} denotes the broader \Aether{} / \Aflow{} framework built on the claims that the \Aether{} is the underlying four-dimensional substrate of reality, the \Aflow{} is its intrinsic ordered motion, observed three-dimensional space is the local experiential slice of that deeper substrate, S-time is the experienced order of change arising from matter, light, and the \Aflow{}, and observed expansion is the three-dimensional appearance of deeper four-dimensional ordered motion. Within that framework, \TheoryName{} denotes the exact-closure theory in which the effective gravitational dynamics are adopted to be exactly Einsteinian with universal matter coupling. In the active sequence, \emph{adoption} means use of that established relativistic dynamics without claiming substrate derivation, while \emph{derivation} is reserved for a first-principles recovery from explicit substrate variables.

The present note stays strictly inside that boundary. It does not claim that final substrate microphysics has already been identified. Its narrower theorem is only that the hidden-translation support reservoir package named by the previous note is itself a canonical exact-shell restriction of a deeper substrate package on the same one-metric line.

\section{Fixed Primitive Continuation Data}

\begin{definition}[Recorded primitive continuation data]
\label{def:fixed}
Fix the following continuation data from the active primitive-reservoir observer-reduction line.
\begin{enumerate}
    \item For each sector label \(\sigma\in\{\Car,\Cur,\Hom\}\), a primitive forcing shell
    \[
    \ForSpace{\sigma},
    \]
    a visible reduction map
    \[
    \Reduction{\sigma}:\ForSpace{\sigma}\to X_{\sigma},
    \]
    and shell-level current and equilibrium carriers
    \[
    \CurrentCarrier{\sigma}:\ForSpace{\sigma}\to Z_{\sigma}^{\mathrm{cur}},
    \qquad
    \EqCarrier{\sigma}:\ForSpace{\sigma}\to Z_{\sigma}^{\mathrm{eq}}.
    \]
    \item The product shell
    \[
    \PrimShell
    \equiv
    \ForSpace{\Car}\times \ForSpace{\Cur}\times \ForSpace{\Hom}.
    \]
    \item Fixed shell extractors
    \[
    \CurrentExtract{\sigma}:Z_{\sigma}^{\mathrm{cur}}\to Y_{\sigma}^{\mathrm{cur}},
    \qquad
    \EqExtract{\sigma}:Z_{\sigma}^{\mathrm{eq}}\to Y_{\sigma}^{\mathrm{eq}},
    \]
    together with the recorded shell composites
    \begin{equation}
    \PrimCurrent{\sigma}
    =
    \CurrentExtract{\sigma}\circ \CurrentCarrier{\sigma},
    \qquad
    \PrimTheta{\sigma}
    =
    \EqExtract{\sigma}\circ \EqCarrier{\sigma}.
    \label{eq:primcomposites}
    \end{equation}
    \item The recorded metric and ordinary-matter outputs
    \[
    \MetricOut:\PrimShell\to\{\text{observer metrics}\},
    \qquad
    \MatterOut:\PrimShell\times\Psi\to\{\text{observer stress tensors}\},
    \]
    satisfying
    \begin{equation}
    \MetricOut(\mathbf w)=\CompletedMetric,
    \qquad
    \MatterOut(\mathbf w,\psi)
    =
    T_{\mu\nu}^{\mathrm{matter}}[\CompletedMetric,\psi]
    \label{eq:fixedoutputs}
    \end{equation}
    for every \(\mathbf w\in\PrimShell\) and every matter configuration \(\psi\in\Psi\).
\end{enumerate}
\end{definition}

\begin{remark}
Definition \ref{def:fixed} is not reopened here. The primitive continuation data, the one operative metric, and the ordinary matter route remain fixed. The new work begins strictly below the previously recorded hidden-translation support reservoir package.
\end{remark}

\section{Recorded Reservoir-Level Target}

The immediately preceding active note records the next burden as the sectorwise \emph{hidden-translation support reservoir package}. In the present manuscript that phrase means exactly the recorded package consisting of:
\begin{enumerate}
    \item a finite-dimensional reservoir state space with compact convex reservoir body, a support target with positive-definite pairing, an exact support section, and normal shell stability;
    \item affine observer-data and shell-response readouts together with the same hidden target and hidden pairing later used at the support layer;
    \item transported symmetry and time-reversal structure;
    \item an exact support shell carrying a zero-core / hidden-seed decomposition and a shell chart into the equation variables;
    \item a linear hidden charge with positive hidden gap;
    \item an observer-compatible exact zero-response realization of the primitive forcing shell;
    \item fixed-data solvability on the support zero core;
    \item the inert zero-core kernel and its orthogonal projector;
    \item and the coercivity inequality on charted data-invisible support-shell directions.
\end{enumerate}
The purpose of the present note is to derive that same recorded package unchanged from a deeper substrate shell-restriction dynamics, and then to re-obtain the same downstream support-layer consequences without changing the benchmark one-metric package.

\section{Primitive Substrate Shell-Restriction Dynamics}

\begin{definition}[Primitive substrate shell-restriction dynamics]
\label{def:substrate}
Fix \(\sigma\in\{\Car,\Cur,\Hom\}\). A \emph{primitive substrate shell-restriction dynamics package} for sector \(\sigma\) consists of the following data.
\begin{enumerate}
    \item A finite-dimensional Euclidean substrate state space
    \[
    \SubSpace{\sigma},
    \]
    a compact convex substrate body
    \[
    \SubBody{\sigma}\subset\SubSpace{\sigma}
    \]
    with nonempty interior, a finite-dimensional Euclidean substrate restriction target
    \[
    \SubTarget{\sigma},
    \]
    a positive-definite symmetric substrate pairing
    \[
    \SubPair{\sigma}:
    \SubTarget{\sigma}\times\SubTarget{\sigma}\to\mathbb R,
    \]
    and a smooth substrate restriction section
    \[
    \SubSection{\sigma}:\SubSpace{\sigma}\to\SubTarget{\sigma}.
    \]
    Define the substrate functional
    \begin{equation}
    \SubFunctional{\sigma}(m)
    \equiv
    \SubPair{\sigma}\bigl(\SubSection{\sigma}(m),\SubSection{\sigma}(m)\bigr)
    \label{eq:subfunctional}
    \end{equation}
    and the exact substrate shell
    \begin{equation}
    \SubShellSet{\sigma}
    \equiv
    \bigl(\SubSection{\sigma}\bigr)^{-1}(0)\cap\SubBody{\sigma}.
    \label{eq:subshell}
    \end{equation}
    Assume \(\SubShellSet{\sigma}\) is a closed embedded submanifold of \(\SubSpace{\sigma}\) and there exists \(\mu_{\sigma}^{\mathrm{sub}}>0\) such that
    \begin{equation}
    \SubFunctional{\sigma}(m)
    \ge
    \mu_{\sigma}^{\mathrm{sub}}\,
    \operatorname{dist}\bigl(m,\SubShellSet{\sigma}\bigr)^{2}
    \qquad
    \text{for all }m\in\SubBody{\sigma}.
    \label{eq:subnormal}
    \end{equation}
    \item Affine substrate observer-data and shell-response readouts
    \[
    \SubData{\sigma}:\SubSpace{\sigma}\to\DataSpace{\sigma},
    \qquad
    \SubShellMap{\sigma}:\SubSpace{\sigma}\to\AmbientTarget{\sigma},
    \]
    a finite-dimensional Euclidean support target
    \[
    \ResTarget{\sigma},
    \]
    a positive-definite symmetric support pairing
    \[
    \ResPair{\sigma}:
    \ResTarget{\sigma}\times\ResTarget{\sigma}\to\mathbb R,
    \]
    a finite-dimensional real hidden target
    \[
    \EqnTarget{\sigma},
    \]
    and a positive-definite symmetric hidden pairing
    \[
    \EqnPair{\sigma}:
    \EqnTarget{\sigma}\times\EqnTarget{\sigma}\to\mathbb R.
    \]
    There is also a compact symmetry group
    \[
    \SubSym{\sigma}
    \]
    and an involution
    \[
    \SubTime{\sigma},
    \]
    acting orthogonally on \(\SubSpace{\sigma}\), \(\DataSpace{\sigma}\), \(\AmbientTarget{\sigma}\), and \(\ResTarget{\sigma}\), and by \(\EqnPair{\sigma}\)-isometries on \(\EqnTarget{\sigma}\), preserving \(\SubBody{\sigma}\), \(\SubShellSet{\sigma}\), \(\SubSection{\sigma}\), \(\SubData{\sigma}\), and \(\SubShellMap{\sigma}\).
    \item A Euclidean shell chart
    \[
    \SubChart{\sigma}:\SubShellSet{\sigma}\to\ResSpace{\sigma}
    \]
    whose shell-body image
    \[
    \ResBody{\sigma}
    \equiv
    \SubChart{\sigma}\bigl(\SubShellSet{\sigma}\bigr)
    \]
    is compact convex with nonempty interior in \(\ResSpace{\sigma}\). There exist affine maps
    \[
    \ResData{\sigma}:\ResSpace{\sigma}\to\DataSpace{\sigma},
    \qquad
    \ResShellMap{\sigma}:\ResSpace{\sigma}\to\AmbientTarget{\sigma},
    \]
    and a smooth support section
    \[
    \ResSection{\sigma}:\ResSpace{\sigma}\to\ResTarget{\sigma}
    \]
    such that
    \begin{equation}
    \ResData{\sigma}\circ\SubChart{\sigma}
    =
    \SubData{\sigma}\big|_{\SubShellSet{\sigma}},
    \qquad
    \ResShellMap{\sigma}\circ\SubChart{\sigma}
    =
    \SubShellMap{\sigma}\big|_{\SubShellSet{\sigma}}.
    \label{eq:subchartreadouts}
    \end{equation}
    There is a smooth shell-internal support recorder
    \[
    \SubSupport{\sigma}:\SubShellSet{\sigma}\to\ResTarget{\sigma}
    \]
    satisfying
    \begin{equation}
    \ResSection{\sigma}\circ\SubChart{\sigma}
    =
    \SubSupport{\sigma}.
    \label{eq:supporttransport}
    \end{equation}
    Define the exact support shell in substrate coordinates by
    \begin{equation}
    \SubSupportShell{\sigma}
    \equiv
    \bigl(\SubSupport{\sigma}\bigr)^{-1}(0)\cap\SubShellSet{\sigma},
    \label{eq:subsupportshell}
    \end{equation}
    and assume
    \begin{equation}
    \ResShellSet{\sigma}
    \equiv
    \bigl(\ResSection{\sigma}\bigr)^{-1}(0)\cap\ResBody{\sigma}
    =
    \SubChart{\sigma}\bigl(\SubSupportShell{\sigma}\bigr).
    \label{eq:transportedsupportshell}
    \end{equation}
    Assume there exists \(\nu_{\sigma}^{\mathrm{sup}}>0\) such that
    \begin{equation}
    \ResPair{\sigma}\bigl(\ResSection{\sigma}(u),\ResSection{\sigma}(u)\bigr)
    \ge
    \nu_{\sigma}^{\mathrm{sup}}\,
    \operatorname{dist}\bigl(u,\ResShellSet{\sigma}\bigr)^{2}
    \qquad
    \text{for all }u\in\ResBody{\sigma}.
    \label{eq:transportednormal}
    \end{equation}
    The symmetry and time-reversal actions transported through \(\SubChart{\sigma}\) act linearly on \(\ResSpace{\sigma}\), preserving \(\ResBody{\sigma}\), \(\ResShellSet{\sigma}\), \(\ResSection{\sigma}\), \(\ResData{\sigma}\), and \(\ResShellMap{\sigma}\).
    \item A closed embedded support zero core
    \[
    \SubZeroCore{\sigma}\subset\SubSupportShell{\sigma},
    \]
    a finite-dimensional Euclidean hidden-seed space
    \[
    \SeedHidden{\sigma},
    \]
    and a compact convex hidden-seed body
    \[
    \SeedHiddenBody{\sigma}\subset\SeedHidden{\sigma}
    \]
    with
    \[
    0\in\operatorname{int}\SeedHiddenBody{\sigma}.
    \]
    There is a Euclidean support-shell chart
    \[
    \ShellChart{\sigma}:\SubSupportShell{\sigma}\to\EqnSpace{\sigma},
    \]
    an affine chart on the support zero core
    \[
    \SubZeroChart{\sigma}:\SubZeroCore{\sigma}\to\EqnZero{\sigma},
    \]
    and a linear isometric embedding
    \[
    \HiddenChart{\sigma}:\SeedHidden{\sigma}\to\EqnHidden{\sigma}
    \]
    such that
    \[
    \EqnSpace{\sigma}
    =
    \EqnZero{\sigma}\oplus\EqnHidden{\sigma}
    \]
    is an orthogonal direct sum, the set
    \[
    \ZeroBody{\sigma}
    \equiv
    \SubZeroChart{\sigma}\bigl(\SubZeroCore{\sigma}\bigr)
    \]
    is compact convex with nonempty interior in \(\EqnZero{\sigma}\), the shell image
    \[
    \EqnBody{\sigma}
    \equiv
    \ShellChart{\sigma}\bigl(\SubSupportShell{\sigma}\bigr)
    \]
    has nonempty interior in \(\EqnSpace{\sigma}\), and there is a diffeomorphism
    \begin{equation}
    \SubSplit{\sigma}:
    \SubZeroCore{\sigma}
    \times
    \SeedHiddenBody{\sigma}
    \xrightarrow{\cong}
    \SubSupportShell{\sigma}
    \label{eq:subsplit}
    \end{equation}
    satisfying
    \begin{equation}
    \ShellChart{\sigma}\bigl(\SubSplit{\sigma}(z,\eta)\bigr)
    =
    \SubZeroChart{\sigma}(z)+\HiddenChart{\sigma}(\eta).
    \label{eq:subchartsplit}
    \end{equation}
    The shell readouts become affine in the support-shell chart: there exist affine maps
    \[
    \EqnData{\sigma}:\EqnSpace{\sigma}\to\DataSpace{\sigma},
    \qquad
    \EqnShell{\sigma}:\EqnSpace{\sigma}\to\AmbientTarget{\sigma},
    \]
    such that
    \begin{equation}
    \EqnData{\sigma}\circ\ShellChart{\sigma}
    =
    \SubData{\sigma}\big|_{\SubSupportShell{\sigma}},
    \qquad
    \EqnShell{\sigma}\circ\ShellChart{\sigma}
    =
    \SubShellMap{\sigma}\big|_{\SubSupportShell{\sigma}}.
    \label{eq:eqnchartreadouts}
    \end{equation}
    The transported symmetry and time-reversal actions act linearly on \(\EqnSpace{\sigma}\), preserving \(\EqnZero{\sigma}\), \(\EqnHidden{\sigma}\), and \(\EqnBody{\sigma}\).
    \item A linear hidden charge map
    \[
    \SubCharge{\sigma}:\SeedHidden{\sigma}\to\EqnTarget{\sigma}
    \]
    such that there exists \(\lambda_{\sigma}^{\mathrm{eqn}}>0\) with
    \begin{equation}
    \EqnPair{\sigma}\bigl(\SubCharge{\sigma}\eta,\SubCharge{\sigma}\eta\bigr)
    \ge
    \lambda_{\sigma}^{\mathrm{eqn}}\,
    \|\eta\|^{2}
    \qquad
    \text{for all }\eta\in\SeedHidden{\sigma}.
    \label{eq:subchargegap}
    \end{equation}
    Define the support-shell hidden recorder by
    \begin{equation}
    \SubResidual{\sigma}\bigl(\SubSplit{\sigma}(z,\eta)\bigr)
    \equiv
    \SubCharge{\sigma}(\eta).
    \label{eq:subresidual}
    \end{equation}
    \item A smooth embedding
    \[
    \jmath_{\sigma}:X_{\sigma}\hookrightarrow\VisibleAffine{\sigma}
    \]
    and a smooth observer-compatible exact zero-response realization
    \[
    \SubEmbed{\sigma}:\ForSpace{\sigma}\hookrightarrow\SubZeroCore{\sigma}
    \]
    whose image is exactly
    \[
    \SubZeroCore{\sigma}
    \cap
    \bigl(\SubShellMap{\sigma}\bigr)^{-1}(0),
    \]
    and such that
    \begin{equation}
    \SubData{\sigma}\bigl(\SubEmbed{\sigma}(w)\bigr)
    =
    \bigl(
    \jmath_{\sigma}(\Reduction{\sigma}(w)),
    \CurrentCarrier{\sigma}(w),
    \EqCarrier{\sigma}(w)
    \bigr)
    \label{eq:subcompat}
    \end{equation}
    for every \(w\in\ForSpace{\sigma}\).
    \item Fixed-data solvability on the support zero core:
    \begin{equation}
    \SubData{\sigma}\bigl(\SubZeroCore{\sigma}\bigr)
    =
    \SubData{\sigma}\Bigl(
    \SubZeroCore{\sigma}\cap
    \bigl(\SubShellMap{\sigma}\bigr)^{-1}(0)
    \Bigr).
    \label{eq:subsolvability}
    \end{equation}
    \item The inert support-zero-core kernel
    \[
    \SubKernel{\sigma}
    \equiv
    \ker\bigl(D\SubData{\sigma}\big|_{\SubZeroCore{\sigma}}\bigr)
    \cap
    \ker\bigl(D\SubShellMap{\sigma}\big|_{\SubZeroCore{\sigma}}\bigr)
    \]
    and the orthogonal projector
    \[
    \SubStateComp{\sigma}:
    \EqnZero{\sigma}\to
    \bigl(\SubKernelCh{\sigma}\bigr)^{\perp}\cap\EqnZero{\sigma},
    \qquad
    \SubKernelCh{\sigma}
    \equiv
    D\SubZeroChart{\sigma}\bigl(\SubKernel{\sigma}\bigr).
    \]
    There exists \(\kappa_{\sigma}^{\mathrm{eqn}}>0\) such that for every
    \[
    w\in\ForSpace{\sigma},
    \qquad
    \delta m\in T_{\SubEmbed{\sigma}(w)}\SubSupportShell{\sigma}
    \]
    satisfying
    \[
    D\SubData{\sigma}\,\delta m=0,
    \]
    if
    \[
    D\ShellChart{\sigma}(\delta m)=\delta z+\delta\eta,
    \qquad
    \delta z\in\EqnZero{\sigma},
    \quad
    \delta\eta\in\EqnHidden{\sigma},
    \]
    then
    \begin{equation}
    \|D\SubShellMap{\sigma}\,\delta m\|^{2}
    +
    \EqnPair{\sigma}\bigl(
    \SubCharge{\sigma}\bigl((\HiddenChart{\sigma})^{-1}\delta\eta\bigr),
    \SubCharge{\sigma}\bigl((\HiddenChart{\sigma})^{-1}\delta\eta\bigr)
    \bigr)
    \ge
    \kappa_{\sigma}^{\mathrm{eqn}}
    \left(
    \|\SubStateComp{\sigma}(\delta z)\|^{2}
    +
    \|\delta\eta\|^{2}
    \right).
    \label{eq:subcoercive}
    \end{equation}
\end{enumerate}
\end{definition}

\begin{remark}
Definition \ref{def:substrate} starts at the correct deeper layer for the present project goal. The reservoir body is no longer primitive there; it is the chart image of the exact substrate shell. The reservoir support section is no longer inserted independently; it is the transport of the shell-internal support recorder. The zero-core / hidden-seed decomposition, shell chart, hidden charge, compatibility realization, solvability statement, inert kernel/projector, and coercivity inequality all already live on the support shell cut out inside the deeper substrate shell. This is therefore an origin-side move below the recorded reservoir package, not another restatement of that package.
\end{remark}

\section{Induced Hidden-Translation Support Reservoir Package}

\begin{proposition}[The recorded reservoir package is produced unchanged]
\label{prop:reservoir}
Assume Definition \ref{def:substrate}. Define, for each sector \(\sigma\in\{\Car,\Cur,\Hom\}\),
\begin{equation}
\ResZeroCore{\sigma}
\equiv
\SubChart{\sigma}\bigl(\SubZeroCore{\sigma}\bigr),
\qquad
\ResEmbed{\sigma}
\equiv
\SubChart{\sigma}\circ\SubEmbed{\sigma},
\label{eq:inducedreszero}
\end{equation}
transport the support-shell chart and zero-core chart through \(\SubChart{\sigma}\), and continue to denote the transported charts by \(\ShellChart{\sigma}\) and \(\ZeroChart{\sigma}\). Concretely,
\begin{equation}
\ZeroChart{\sigma}
\equiv
\SubZeroChart{\sigma}\circ
\bigl(\SubChart{\sigma}\big|_{\SubZeroCore{\sigma}}\bigr)^{-1},
\label{eq:transportedzerochart}
\end{equation}
while on the reservoir support shell the transported shell chart is the map
\[
u\longmapsto
\ShellChart{\sigma}\Bigl(
\bigl(\SubChart{\sigma}\big|_{\SubSupportShell{\sigma}}\bigr)^{-1}(u)
\Bigr),
\]
again denoted by \(\ShellChart{\sigma}\).
Define the transported split map on the reservoir support shell by
\begin{equation}
\SplitMap{\sigma}(r,\eta)
\equiv
\SubChart{\sigma}\Bigl(
\SubSplit{\sigma}\bigl(
(\SubChart{\sigma}\big|_{\SubZeroCore{\sigma}})^{-1}(r),
\eta
\bigr)
\Bigr).
\label{eq:transportedsplit}
\end{equation}
Then the resulting data satisfy exactly the recorded hidden-translation support reservoir package of the preceding theorem note, with the same reservoir body \(\ResBody{\sigma}\), the same support section \(\ResSection{\sigma}\), the same readouts \(\ResData{\sigma},\ResShellMap{\sigma}\), the same hidden pairing \(\EqnPair{\sigma}\), the same hidden-seed space \(\SeedHidden{\sigma}\), and the same hidden charge map
\[
\HiddenCharge{\sigma}\equiv\SubCharge{\sigma}.
\]
\end{proposition}

\begin{proof}
The reservoir state space \(\ResSpace{\sigma}\), the reservoir body \(\ResBody{\sigma}\), the support target \(\ResTarget{\sigma}\), the support pairing \(\ResPair{\sigma}\), the support section \(\ResSection{\sigma}\), and the exact support shell \(\ResShellSet{\sigma}\) are already part of Definition \ref{def:substrate}(2)--(3). Equation \eqref{eq:transportednormal} is exactly the required normal shell stability on the reservoir body.

The affine reservoir readouts \(\ResData{\sigma}\) and \(\ResShellMap{\sigma}\) are given by Definition \ref{def:substrate}(3), and \eqref{eq:subchartreadouts} identifies them as the chart transports of the deeper substrate readouts. The transported symmetry and time-reversal actions on \(\ResSpace{\sigma}\) are also part of Definition \ref{def:substrate}(3), and they preserve the reservoir body, support shell, support section, and readouts exactly as required.

Definition \ref{def:substrate}(4) provides the support zero core, hidden-seed space, hidden-seed body, support-shell chart, zero-core chart, and hidden chart. Transporting the support-shell chart as prescribed above and the zero-core chart by \eqref{eq:transportedzerochart} gives the reservoir-shell chart and reservoir zero-core chart. The transported split map \eqref{eq:transportedsplit} is a diffeomorphism because \(\SubSplit{\sigma}\) is one and \(\SubChart{\sigma}\) is a shell chart. Moreover,
\[
\ShellChart{\sigma}\bigl(\SplitMap{\sigma}(r,\eta)\bigr)
=
\ZeroChart{\sigma}(r)+\HiddenChart{\sigma}(\eta),
\]
so the recorded reservoir shell decomposition is recovered unchanged.

The hidden charge is exactly \(\HiddenCharge{\sigma}=\SubCharge{\sigma}\), and the positive hidden gap is \eqref{eq:subchargegap}. Transporting \eqref{eq:subresidual} through \(\SubChart{\sigma}\) gives the reservoir hidden recorder
\[
\ResResidual{\sigma}\bigl(\SplitMap{\sigma}(r,\eta)\bigr)
=
\HiddenCharge{\sigma}(\eta),
\]
which is precisely the recorded reservoir residual.

The observer-compatible exact zero-response realization is \(\ResEmbed{\sigma}=\SubChart{\sigma}\circ\SubEmbed{\sigma}\). Because \(\SubEmbed{\sigma}\) lands in \(\SubZeroCore{\sigma}\cap(\SubShellMap{\sigma})^{-1}(0)\), the transported map lands in \(\ResZeroCore{\sigma}\cap(\ResShellMap{\sigma})^{-1}(0)\). Equation \eqref{eq:subcompat} and the chart transport \eqref{eq:subchartreadouts} give the required compatibility identity on the reservoir side.

The solvability statement is the chart image of \eqref{eq:subsolvability}. Likewise the inert kernel and orthogonal projector transport from \(\SubZeroCore{\sigma}\) to \(\ResZeroCore{\sigma}\) via the zero-core chart. Finally, the chart transport of \eqref{eq:subcoercive} gives the recorded reservoir coercivity inequality on data-invisible directions tangent to the reservoir support shell. Every item of the recorded hidden-translation support reservoir package is therefore recovered unchanged from the deeper substrate shell-restriction dynamics.
\end{proof}

\begin{remark}
Proposition \ref{prop:reservoir} is the key gain of the present note. The hidden-translation support reservoir package is no longer a reservoir-level insertion on the active line. It is the exact-shell restriction of a more primitive substrate dynamics, with the same body, the same section, the same zero-core / hidden-seed split, the same hidden charge, and the same coercive support-shell control.
\end{remark}

\section{Induced Support-Layer Consequences}

\begin{proposition}[The same designated support package and hidden-fiber exactness are forced]
\label{prop:support}
Assume Definition \ref{def:substrate}. Define the common zero-locus defect map on the orthogonal decomposition
\[
\EqnSpace{\sigma}=\EqnZero{\sigma}\oplus\EqnHidden{\sigma}
\]
by
\begin{equation}
\EqnLin{\sigma}(z+\eta)
\equiv
\SubCharge{\sigma}\bigl((\HiddenChart{\sigma})^{-1}\eta\bigr),
\qquad
z\in\EqnZero{\sigma},
\ \eta\in\EqnHidden{\sigma},
\label{eq:inducedlin}
\end{equation}
and define the residual
\begin{equation}
\EqnResidual{\sigma}(z+\eta)
\equiv
\EqnLin{\sigma}(\eta).
\label{eq:inducedresidual}
\end{equation}
Then the designated support-layer data and hidden-fiber exactness of the preceding theorem note are recovered unchanged. More precisely:
\begin{enumerate}
    \item the equation-state body is
    \[
    \EqnBody{\sigma}
    =
    \left\{
    z+\eta:
    z\in\ZeroBody{\sigma},
    \ \eta\in\HiddenChart{\sigma}\bigl(\SeedHiddenBody{\sigma}\bigr)
    \right\},
    \]
    hence compact convex with nonempty interior in \(\EqnSpace{\sigma}\);
    \item the readouts \(\EqnData{\sigma}\) and \(\EqnShell{\sigma}\) are affine, the symmetry and time-reversal actions on \(\EqnSpace{\sigma}\) are linear and preserve \(\EqnZero{\sigma}\), \(\EqnHidden{\sigma}\), and \(\EqnBody{\sigma}\), and \(\EqnLin{\sigma}\) vanishes on \(\EqnZero{\sigma}\) while obeying the positive hidden gap
    \begin{equation}
    \EqnPair{\sigma}\bigl(\EqnLin{\sigma}\eta,\EqnLin{\sigma}\eta\bigr)
    \ge
    \lambda_{\sigma}^{\mathrm{eqn}}\,
    \|\eta\|^{2}
    \qquad
    \text{for all }\eta\in\EqnHidden{\sigma};
    \label{eq:hidgap}
    \end{equation}
    \item the exact zero-response realization
    \[
    \EqnEmbed{\sigma}
    \equiv
    \ShellChart{\sigma}\circ\SubEmbed{\sigma}
    \]
    lands in
    \[
    \EqnZero{\sigma}\cap\EqnBody{\sigma}\cap\bigl(\EqnShell{\sigma}\bigr)^{-1}(0)
    \]
    and satisfies
    \begin{equation}
    \EqnData{\sigma}\bigl(\EqnEmbed{\sigma}(w)\bigr)
    =
    \bigl(
    \jmath_{\sigma}(\Reduction{\sigma}(w)),
    \CurrentCarrier{\sigma}(w),
    \EqCarrier{\sigma}(w)
    \bigr)
    \label{eq:eqncompat}
    \end{equation}
    for every \(w\in\ForSpace{\sigma}\);
    \item fixed-data solvability, the inert kernel/projector, and the coercivity inequality
    \begin{equation}
    \|D\EqnShell{\sigma}\,\delta u\|^{2}
    +
    \EqnPair{\sigma}\bigl(\EqnLin{\sigma}\delta u,\EqnLin{\sigma}\delta u\bigr)
    \ge
    \kappa_{\sigma}^{\mathrm{eqn}}
    \left(
    \|\StateComp{\sigma}\bigl((I-\EqnProj{\sigma})\delta u\bigr)\|^{2}
    +
    \|\EqnProj{\sigma}\delta u\|^{2}
    \right)
    \label{eq:eqncoercive}
    \end{equation}
    hold on data-invisible directions at \(\EqnEmbed{\sigma}(w)\);
    \item the residual \(\EqnResidual{\sigma}\) is zero-fiber invariant, hidden additive, support-linearization compatible, and on the support shell agrees with the deeper hidden recorder:
    \begin{equation}
    \EqnResidual{\sigma}\bigl(\ShellChart{\sigma}(\SubSplit{\sigma}(z,\eta))\bigr)
    =
    \SubResidual{\sigma}\bigl(\SubSplit{\sigma}(z,\eta)\bigr).
    \label{eq:residualagreement}
    \end{equation}
\end{enumerate}
\end{proposition}

\begin{proof}
Item (1) is immediate from \eqref{eq:subsplit}--\eqref{eq:subchartsplit}. The shell body is the image of the product of two compact convex bodies under the affine addition map \((z,\eta)\mapsto z+\eta\), so it is compact convex. It has nonempty interior because \(\ZeroBody{\sigma}\) has nonempty interior in \(\EqnZero{\sigma}\) and \(0\in\operatorname{int}\SeedHiddenBody{\sigma}\) while \(\HiddenChart{\sigma}\) is linear isometric.

Item (2) is built into Definition \ref{def:substrate}(4) for the affine readouts and linear transported symmetries. The positive hidden gap \eqref{eq:hidgap} is the charted form of \eqref{eq:subchargegap} because \(\HiddenChart{\sigma}\) is isometric. By construction \(\EqnLin{\sigma}\) depends only on the hidden component, so it vanishes on \(\EqnZero{\sigma}\).

For item (3), \(\SubEmbed{\sigma}\) lands in the support zero core, hence \(\EqnEmbed{\sigma}(\ForSpace{\sigma})\subset\EqnZero{\sigma}\). Equation \eqref{eq:eqncompat} follows by combining \eqref{eq:subcompat} with \eqref{eq:eqnchartreadouts}. Because \(\SubEmbed{\sigma}\bigl(\ForSpace{\sigma}\bigr)=\SubZeroCore{\sigma}\cap(\SubShellMap{\sigma})^{-1}(0)\), its chart image is exactly the displayed zero-response locus in the equation variables.

For item (4), fixed-data solvability is the chart image of \eqref{eq:subsolvability}. Define
\[
\Kernel{\sigma}
\equiv
\SubKernelCh{\sigma},
\qquad
\StateComp{\sigma}
\equiv
\SubStateComp{\sigma},
\qquad
\EqnProj{\sigma}:\EqnSpace{\sigma}\to\EqnHidden{\sigma}
\]
to be the orthogonal projector onto the hidden summand. Let \(w\in\ForSpace{\sigma}\) and \(\delta u\in\EqnSpace{\sigma}\) satisfy \(D\EqnData{\sigma}\,\delta u=0\). Write \(\delta u=\delta z+\delta\eta\) with \(\delta z\in\EqnZero{\sigma}\) and \(\delta\eta\in\EqnHidden{\sigma}\), and define
\[
\delta m
\equiv
D\bigl((\ShellChart{\sigma})^{-1}\bigr)(\delta u)
\in
T_{\SubEmbed{\sigma}(w)}\SubSupportShell{\sigma}.
\]
Then \(D\SubData{\sigma}\,\delta m=0\), so \eqref{eq:subcoercive} yields
\[
\|D\EqnShell{\sigma}\,\delta u\|^{2}
+
\EqnPair{\sigma}\bigl(\EqnLin{\sigma}\delta u,\EqnLin{\sigma}\delta u\bigr)
\ge
\kappa_{\sigma}^{\mathrm{eqn}}
\left(
\|\StateComp{\sigma}(\delta z)\|^{2}
+
\|\delta\eta\|^{2}
\right).
\]
Since \(\delta z=(I-\EqnProj{\sigma})\delta u\) and \(\delta\eta=\EqnProj{\sigma}\delta u\), this is exactly \eqref{eq:eqncoercive}.

For item (5), the agreement formula \eqref{eq:residualagreement} follows from \eqref{eq:subchartsplit}, \eqref{eq:subresidual}, and \eqref{eq:inducedlin}:
\[
\EqnResidual{\sigma}\bigl(\ShellChart{\sigma}(\SubSplit{\sigma}(z,\eta))\bigr)
=
\EqnResidual{\sigma}\bigl(\SubZeroChart{\sigma}(z)+\HiddenChart{\sigma}(\eta)\bigr)
=
\EqnLin{\sigma}\bigl(\HiddenChart{\sigma}(\eta)\bigr)
=
\SubCharge{\sigma}(\eta).
\]
If \(u=z+\eta\) with \(z\in\EqnZero{\sigma}\) and \(\eta\in\EqnHidden{\sigma}\), then
\[
\EqnResidual{\sigma}(u)
=
\EqnLin{\sigma}(\eta)
=
\EqnLin{\sigma}\bigl(\EqnProj{\sigma}u\bigr)
=
\EqnResidual{\sigma}\bigl(\EqnProj{\sigma}u\bigr),
\]
which is zero-fiber invariance. If \(\eta_{1},\eta_{2}\in\EqnHidden{\sigma}\), linearity of \(\EqnLin{\sigma}\) gives
\[
\EqnResidual{\sigma}(\eta_{1}+\eta_{2})
=
\EqnResidual{\sigma}(\eta_{1})
+
\EqnResidual{\sigma}(\eta_{2}),
\]
which is hidden additivity. Finally, \(\EqnResidual{\sigma}\) is linear and agrees with \(\EqnLin{\sigma}\) on all of \(\EqnSpace{\sigma}\), so for every \(z\in\EqnZero{\sigma}\),
\[
D\EqnResidual{\sigma}(z)=\EqnLin{\sigma}.
\]
This is the required support-linearization compatibility.
\end{proof}

\begin{lemma}[The induced residual is the canonical hidden linear defect]
\label{lem:canonical}
Assume the hypotheses of Proposition \ref{prop:support}. Then for every \(u\in\EqnSpace{\sigma}\),
\begin{equation}
\EqnResidual{\sigma}(u)
=
\EqnLin{\sigma}\bigl(\EqnProj{\sigma}u\bigr)
=
\EqnLin{\sigma}(u).
\label{eq:residualcanonical}
\end{equation}
\end{lemma}

\begin{proof}
Write \(u=z+\eta\) with \(z\in\EqnZero{\sigma}\) and \(\eta=\EqnProj{\sigma}u\in\EqnHidden{\sigma}\). Then
\[
\EqnResidual{\sigma}(u)
=
\EqnLin{\sigma}(\eta)
=
\EqnLin{\sigma}\bigl(\EqnProj{\sigma}u\bigr).
\]
Since \(\EqnLin{\sigma}\) vanishes on \(\EqnZero{\sigma}\),
\[
\EqnLin{\sigma}(u)
=
\EqnLin{\sigma}(z+\eta)
=
\EqnLin{\sigma}(\eta),
\]
and \eqref{eq:residualcanonical} follows.
\end{proof}

\begin{remark}
The downstream gain is exactly the one the project needs. Starting below the reservoir package does not change the support-layer output; it explains why that output exists. The same support data and the same hidden-fiber exactness are now charted consequences of a deeper substrate shell-restriction package.
\end{remark}

\section{Benchmark Preservation}

\begin{theorem}[The same benchmark one-metric package is preserved]
\label{thm:outputs}
Assume Definition \ref{def:fixed} and, for each sector \(\sigma\in\{\Car,\Cur,\Hom\}\), a primitive substrate shell-restriction dynamics package in the sense of Definition \ref{def:substrate}. Then the benchmark output statements of \eqref{eq:fixedoutputs} remain unchanged:
\[
\MetricOut(\mathbf w)=\CompletedMetric,
\qquad
\MatterOut(\mathbf w,\psi)
=
T_{\mu\nu}^{\mathrm{matter}}[\CompletedMetric,\psi]
\]
for every \(\mathbf w\in\PrimShell\) and every matter configuration \(\psi\in\Psi\). Consequently the deeper substrate shell-restriction theorem introduces neither a second operative metric nor an enlarged low-energy matter law.
\end{theorem}

\begin{proof}
The present note changes only the deeper origin of the reservoir and support-layer data. It does not alter the fixed primitive shell \(\PrimShell\), the recorded metric map \(\MetricOut\), or the recorded matter map \(\MatterOut\). Those remain the continuation data of Definition \ref{def:fixed}, so \eqref{eq:fixedoutputs} continues to hold exactly.
\end{proof}

\section{Main Theorem}

\begin{theorem}[The reservoir-level burden moves below the recorded reservoir package]
\label{thm:main}
Assume the fixed primitive continuation data of Definition \ref{def:fixed} and, for each sector \(\sigma\in\{\Car,\Cur,\Hom\}\), a primitive substrate shell-restriction dynamics package in the sense of Definition \ref{def:substrate}. Then:
\begin{enumerate}
    \item the recorded hidden-translation support reservoir package is produced unchanged by exact shell restriction and chart transport from the deeper substrate dynamics;
    \item the same designated support-layer data are forced again: the compact convex equation-state body, affine observer-data and shell-response readouts, zero/hidden splitting, common zero-locus defect map, exact zero-response realization, fixed-data solvability, inert kernel/projector, and coercivity inequality all follow from the deeper substrate package rather than standing as independent assumptions;
    \item the induced residual is hidden-fiber exact and is exactly the canonical hidden linear defect
    \[
    u\mapsto \EqnLin{\sigma}\bigl(\EqnProj{\sigma}u\bigr),
    \]
    so no new support-level freedom is reintroduced by passing to the deeper layer;
    \item the same benchmark one-metric package remains fixed: the operative metric is still exactly \(\CompletedMetric\), the ordinary matter route is unchanged, there is no second operative metric, and the low-energy matter law is not enlarged;
    \item consequently the remaining honest burden on the active line is now deeper again: derive or justify the primitive substrate shell-restriction dynamics package itself from still more primitive substrate dynamics, or else prove a sharper inevitability theorem at that deeper shell-restriction level.
\end{enumerate}
\end{theorem}

\begin{proof}
Item (1) is Proposition \ref{prop:reservoir}. Items (2) and (3) are Proposition \ref{prop:support} together with Lemma \ref{lem:canonical}. Item (4) is Theorem \ref{thm:outputs}. Item (5) is the routing consequence of the previous items.
\end{proof}

\begin{corollary}[What remains open after this theorem]
\label{cor:next}
After Theorem \ref{thm:main}, the remaining honest burden is no longer to explain the hidden-translation support reservoir package itself. That package is now the exact-shell restriction of a deeper substrate dynamics on the active one-metric line. What remains open is why the primitive substrate shell-restriction package should hold, or whether an even sharper inevitability theorem can remove that still deeper freedom while preserving the same benchmark one-metric package and claim boundary.
\end{corollary}

\begin{proof}
Theorem \ref{thm:main} identifies the recorded reservoir package as a derived substrate-shell consequence rather than the new disciplined stopping point. Therefore the next honest burden lies below that package.
\end{proof}

\section{Conclusion}

The positive result of this note is narrower than a completed first-principles derivation and stronger than another reservoir-level restatement. The hidden-translation support reservoir package is no longer left standing as the next unexplained insertion on the active line. Starting from a deeper primitive substrate shell-restriction dynamics, the note derives the reservoir body as the chart image of an exact substrate shell, derives the reservoir support section as the transported shell-internal support recorder, preserves the same normal stability, and transports the same zero-core / hidden-seed decomposition, shell chart, hidden charge, compatibility realization, fixed-data solvability, inert kernel/projector, and coercive control.

Because the same reservoir package is recovered unchanged, the same downstream support-layer consequence survives unchanged as well. The designated equation-state body, affine readouts, zero/hidden splitting, common defect map, exact zero-response realization, solvability statement, inert kernel/projector, and coercivity inequality are all forced again. The residual remains exactly the canonical hidden linear defect \(u\mapsto\EqnLin(\EqnProj u)\), so hidden-fiber exactness is retained without introducing any new support-level freedom.

The scope remains disciplined. The benchmark package is unchanged: the one operative metric is still exactly \(\CompletedMetric\), the ordinary matter route is unchanged, there is no second operative metric, and the low-energy matter law is not enlarged. This is therefore an origin-side derivational gain inside the active program of \TheoryName{}, not a basis for stronger promotion language. The next honest burden is deeper again: derive or justify the primitive substrate shell-restriction dynamics package itself from still more primitive substrate structure, or else prove a sharper inevitability theorem there. Until then, adoption and derivation should continue to be kept sharply separated.

\end{document}
